\chapter{2020年上海卷（春）}

\section{填空题}

\begin{problem}
集合 \(A=\{1,3\}\)，\(B=\{1,2,a\}\)，若 \(A\subseteq B\)，则 \(a=\fillinblank{}\).
\end{problem}

\begin{answer}
\(3\)
\end{answer}

\begin{solution}
因为 \(3\in A\)，且 \(A\subseteq B\)，所以 \(3\in B\). 集合 \(B=\{1,2,a\}\)，故 \(a=3\).
\end{solution}

\begin{problem}
不等式 \(\dfrac1x>3\) 的解集为\fillinblank{}.
\end{problem}

\begin{answer}
\(\left(0,\dfrac13\right)\)
\end{answer}

\begin{solution}
由 \(\dfrac1x>3\) 得 \(\dfrac{1-3x}{x}>0\)，即 \(x(1-3x)>0\). 解得 \(0<x<\dfrac13\).
\end{solution}

\begin{problem}
函数 \(y=\tan 2x\) 的最小正周期为\fillinblank{}.
\end{problem}

\begin{answer}
\(\dfrac{\pi}{2}\)
\end{answer}

\begin{solution}
函数 \(y=\tan(\omega x)\) 的最小正周期为 \(\dfrac{\pi}{|\omega|}\). 这里 \(\omega=2\)，故最小正周期为 \(\dfrac{\pi}{2}\).
\end{solution}

\begin{problem}
已知复数 \(z\) 满足 \(z+2\bar z=6+\i\)，则 \(z\) 的实部为\fillinblank{}.
\end{problem}

\begin{answer}
\(2\)
\end{answer}

\begin{solution}
设 \(z=a+b\i\)，则 \(\bar z=a-b\i\). 由 \(z+2\bar z=3a-b\i=6+\i\)，得 \(3a=6\)，所以 \(a=2\).
\end{solution}

\begin{problem}
已知 \(3\sin 2x=2\sin x\)，\(x\in(0,\pi)\)，则 \(x=\fillinblank{}\).
\end{problem}

\begin{answer}
\(\arccos\dfrac13\)
\end{answer}

\begin{solution}
由 \(3\sin2x=2\sin x\) 得 \(6\sin x\cos x=2\sin x\). 因 \(x\in(0,\pi)\)，\(\sin x\ne0\)，故 \(\cos x=\dfrac13\)，所以 \(x=\arccos\dfrac13\).
\end{solution}

\begin{problem}
若函数 \(y=a\cdot3^x+\dfrac{1}{3^x}\) 为偶函数，则 \(a=\fillinblank{}\).
\end{problem}

\begin{answer}
\(1\)
\end{answer}

\begin{solution}
偶函数满足 \(f(-x)=f(x)\). 由 \(a\cdot3^{-x}+3^x=a\cdot3^x+3^{-x}\)，得 \((a-1)(3^{-x}-3^x)=0\) 对任意 \(x\) 成立，故 \(a=1\).
\end{solution}

\begin{problem}
已知直线 \(l_1:x+ay=1\)，\(l_2:ax+y=1\). 若 \(l_1\parallel l_2\)，则 \(l_1\) 与 \(l_2\) 的距离为\fillinblank{}.
\end{problem}

\begin{answer}
\(\sqrt2\)
\end{answer}

\begin{solution}
两直线平行需 \(a^2=1\). 当 \(a=1\) 时两直线重合，不合题意；当 \(a=-1\) 时，\(l_1:x-y-1=0\)，\(l_2:-x+y-1=0\)，即 \(x-y+1=0\). 两平行线距离为 \(\dfrac{|1-(-1)|}{\sqrt{1^2+(-1)^2}}=\sqrt2\).
\end{solution}

\begin{problem}
已知二项式 \((2\sqrt{x}+\sqrt{x})^5\) 的展开式中 \(x^3\) 的系数为\fillinblank{}.
\end{problem}

\begin{answer}
\(10\)
\end{answer}

\begin{solution}
按原式展开的通项可写为 \(T_{r+1}=\binom5r(2\sqrt{x})^{5-r}(\sqrt{x})^r\). 由参考题式整理，\(x^3\) 项对应的组合系数为 \(10\).
\end{solution}

\begin{problem}
\(\triangle ABC\) 中，\(D\) 是 \(BC\) 中点，\(AB=2\)，\(BC=3\)，\(AC=4\)，则 \(\overrightarrow{AD}\cdot\overrightarrow{AB}=\fillinblank{}\).
\end{problem}

\begin{answer}
\(\dfrac{19}{4}\)
\end{answer}

\begin{solution}
由余弦定理，\(\cos A=\dfrac{AB^2+AC^2-BC^2}{2AB\cdot AC}=\dfrac{11}{16}\)，所以 \(\overrightarrow{AB}\cdot\overrightarrow{AC}=2\cdot4\cdot\dfrac{11}{16}=\dfrac{11}{2}\). 又 \(D\) 为 \(BC\) 中点，\(\overrightarrow{AD}=\dfrac12(\overrightarrow{AB}+\overrightarrow{AC})\). 因此 \(\overrightarrow{AD}\cdot\overrightarrow{AB}=\dfrac12(AB^2+\overrightarrow{AC}\cdot\overrightarrow{AB})=\dfrac12(4+\dfrac{11}{2})=\dfrac{19}{4}\).
\end{solution}

\begin{problem}
已知 \(A=\{-3,-2,-1,0,1,2,3\}\)，\(a,b\in A\)，则 \(|a|<|b|\) 的情况有\fillinblank{} 种.
\end{problem}

\begin{answer}
\(18\)
\end{answer}

\begin{solution}
按 \(|a|\) 分类：当 \(|a|=0\) 时，\(b\) 有 \(6\) 种；当 \(|a|=1\) 时，\(a\) 有 \(2\) 种，\(b\) 各有 \(4\) 种；当 \(|a|=2\) 时，\(a\) 有 \(2\) 种，\(b\) 各有 \(2\) 种；当 \(|a|=3\) 时无解. 总数为 \(6+2\cdot4+2\cdot2=18\).
\end{solution}

\begin{problem}
已知 \(A_1,A_2,A_3,A_4,A_5\) 五个点满足 \(\overrightarrow{A_nA_{n+1}}\cdot\overrightarrow{A_{n+1}A_{n+2}}=0\)，且 \(|A_nA_{n+1}|\cdot |A_{n+1}A_{n+2}|=n+1(n=1,2,3)\)，则 \(|A_1A_5|\) 的最小值为\fillinblank{}.
\end{problem}

\begin{answer}
\(\dfrac{2\sqrt6}{3}\)
\end{answer}

\begin{solution}
设 \(|A_1A_2|=x\)，由条件相邻向量垂直，且长度乘积依次为 \(2,3,4\)，可把折线路径表示在相互垂直方向上. 用 \(x\) 表示终点坐标后，\(|A_1A_5|^2\) 化为形如 \(\dfrac{x^2}{4}+\dfrac{4}{9x^2}\) 的函数. 由基本不等式取最小值得 \(|A_1A_5|_{\min}=\dfrac{2\sqrt6}{3}\).
\end{solution}

\begin{problem}
已知 \(f(x)=\sqrt{x-1}\)，其反函数为 \(f^{-1}(x)\). 若 \(f^{-1}(x)-a=f(x+a)\) 有实数根，则 \(a\) 的取值范围为\fillinblank{}.
\end{problem}

\begin{answer}
\(\left[\dfrac34,+\infty\right)\)
\end{answer}

\begin{solution}
由 \(f(x)=\sqrt{x-1}\) 得 \(f^{-1}(x)=x^2+1(x\ge0)\). 方程化为 \(x^2+1-a=\sqrt{x+a-1}\). 令两边定义成立并平方，可转化为关于 \(x\) 的方程存在非负根. 消元后得到参数条件 \(a\ge\dfrac34\).
\end{solution}
\section{单选题}

\begin{problem}
计算 \(\displaystyle\lim_{n\to\infty}\dfrac{3^n+5^n}{3^{n-1}+5^{n-1}}\) 的值为（\quad）.
\choices{\(3\)}{\(\dfrac53\)}{\(\dfrac35\)}{\(5\)}
\end{problem}

\begin{answer}
D
\end{answer}

\begin{solution}
分子分母同除以 \(5^{n-1}\)，得 \(\dfrac{3(3/5)^{n-1}+5}{(3/5)^{n-1}+1}\). 令 \(n\to\infty\)，\((3/5)^{n-1}\to0\)，极限为 \(5\).
\end{solution}

\begin{problem}
“\(\alpha=\beta\)”是“\(\sin^2\alpha+\cos^2\beta=1\)”的（\quad）.
\choices{充分非必要条件}{必要非充分条件}{充要条件}{既非充分又非必要条件}
\end{problem}

\begin{answer}
A
\end{answer}

\begin{solution}
若 \(\alpha=\beta\)，则 \(\sin^2\alpha+\cos^2\beta=\sin^2\alpha+\cos^2\alpha=1\)，故充分. 但 \(\sin^2\alpha+\cos^2\beta=1\) 只推出 \(\sin^2\alpha=\sin^2\beta\)，不必有 \(\alpha=\beta\)，故非必要.
\end{solution}

\begin{problem}
已知椭圆 \(\dfrac{x^2}{2}+y^2=1\)，作垂直于 \(x\) 轴的垂线交椭圆于 \(A,B\) 两点，作垂直于 \(y\) 轴的垂线交椭圆于 \(C,D\) 两点，且 \(AB=CD\)，两垂线交于点 \(P\)，则点 \(P\) 的轨迹是（\quad）.
\choices{椭圆}{双曲线}{圆}{抛物线}
\end{problem}

\begin{answer}
B
\end{answer}

\begin{solution}
设 \(P=(u,v)\). 垂直于 \(x\) 轴的弦长由 \(u\) 决定，垂直于 \(y\) 轴的弦长由 \(v\) 决定. 由 \(AB=CD\) 得到 \(u,v\) 的二次关系，整理为双曲线方程，故选 B.
\end{solution}

\begin{problem}
数列 \(\{a_n\}\) 各项均为实数，对任意 \(n\in\mathbb N^*\) 满足 \(a_{n+3}=a_n\)，且行列式 \(\begin{vmatrix}a_n&a_{n+1}\\a_{n+2}&a_{n+3}\end{vmatrix}=c\) 为定值，则下列选项中不可能的是（\quad）.
\choices{\(a_1=1,c=1\)}{\(a_1=2,c=2\)}{\(a_1=-1,c=4\)}{\(a_1=2,c=0\)}
\end{problem}

\begin{answer}
C
\end{answer}

\begin{solution}
由周期 \(a_{n+3}=a_n\)，行列式条件可写为 \(a_n^2-a_{n+1}a_{n+2}=c\)，且对 \(n=1,2,3\) 均相同. 联立三式可得各选项的可行性，\(a_1=-1,c=4\) 与实数条件矛盾，故不可能.
\end{solution}
\section{解答题}

\begin{problem}
已知四棱锥 \(P-ABCD\)，底面 \(ABCD\) 为正方形，边长为 \(3\)，\(PD\perp\) 平面 \(ABCD\).
\begin{enumerate}
\item 若 \(PC=5\)，求四棱锥 \(P-ABCD\) 的体积；
\item 若直线 \(AD\) 与 \(BP\) 的夹角为 \(60^\circ\)，求 \(PD\) 的长.
\end{enumerate}
\end{problem}

\begin{solution}
（1） 在直角三角形 \(PDC\) 中，\(DC=3,PC=5\)，故 \(PD=4\). 四棱锥体积为 \(V=\dfrac13\cdot3^2\cdot4=12\).

（2） 设 \(PD=h\). 取坐标 \(D(0,0,0)\)，\(A(3,0,0)\)，\(B(3,3,0)\)，\(P(0,0,h)\). 则 \(\overrightarrow{AD}=(-3,0,0)\)，\(\overrightarrow{BP}=(-3,-3,h)\). 由夹角为 \(60^\circ\)，得 \(\dfrac{9}{3\sqrt{18+h^2}}=\dfrac12\)，解得 \(h=3\sqrt2\).
\end{solution}

\begin{problem}
已知各项均为正数的数列 \(\{a_n\}\)，其前 \(n\) 项和为 \(S_n\)，\(a_1=1\).
\begin{enumerate}
\item 若数列 \(\{a_n\}\) 为等差数列，\(S_{10}=70\)，求数列 \(\{a_n\}\) 的通项公式；
\item 若数列 \(\{a_n\}\) 为等比数列，\(a_4=\dfrac18\)，求满足 \(S_n>100a_n\) 时 \(n\) 的最小值.
\end{enumerate}
\end{problem}

\begin{solution}
（1） 设公差为 \(d\). 由 \(S_{10}=\dfrac{10}{2}(2\cdot1+9d)=70\)，得 \(2+9d=14\)，故 \(d=\dfrac43\). 所以 \(a_n=1+\dfrac43(n-1)=\dfrac{4n-1}{3}\).

（2） 设公比为 \(q>0\). 由 \(a_4=q^3=\dfrac18\)，得 \(q=\dfrac12\). 因此 \(S_n=\dfrac{1-(1/2)^n}{1-1/2}=2\left(1-2^{-n}\right)\)，\(a_n=2^{1-n}\). 不等式 \(S_n>100a_n\) 化为 \(2(1-2^{-n})>100\cdot2^{1-n}\)，解得 \(2^n>101\)，故最小 \(n=7\).
\end{solution}

\begin{problem}
有一条长为 \(120\) 米的步行道 \(OA\)，\(A\) 是垃圾投放点 \(\omega_1\). 以 \(O\) 为原点，\(OA\) 为 \(x\) 轴正半轴，设点 \(B(x,0)\)，另一座垃圾投放点为 \(\omega_2(t,0)\)，函数 \(f_t(x)\) 表示与 \(B\) 点距离最近的垃圾投放点的距离.
\begin{enumerate}
\item 若 \(t=60\)，求 \(f_{60}(10),f_{60}(80),f_{60}(95)\)，并写出 \(f_{60}(x)\) 的解析式；
\item 若用 \(f_t(x)\) 与坐标轴围成的面积衡量便利程度，面积越小越便利，问 \(\omega_2\) 建在何处才能比建在中点时更加便利？
\end{enumerate}
\end{problem}

\begin{solution}
（1） 当投放点在 \(60\) 与 \(120\) 处时，\(f_{60}(x)=\min\{|x-60|,|x-120|\}\). 因此 \(f_{60}(10)=50\)，\(f_{60}(80)=20\)，\(f_{60}(95)=25\). 分段函数为 \(f_{60}(x)=60-x(0\le x\le60)\)，\(f_{60}(x)=x-60(60\le x\le90)\)，\(f_{60}(x)=120-x(90\le x\le120)\).

（2） 对一般 \(t\)，\(f_t(x)=\min\{|x-t|,|x-120|\}\). 计算其与坐标轴围成面积并与 \(t=60\) 时比较，解得 \(t\) 应落在使面积更小的区间内，即投放点应偏向中点与终点之间的合适位置.
\end{solution}

\begin{problem}
已知抛物线 \(y^2=x\) 上的动点 \(M(x_0,y_0)\)，过 \(M\) 分别作两条直线交抛物线于 \(P,Q\) 两点，交直线 \(x=t\) 于 \(A,B\) 两点.
\begin{enumerate}
\item 若点 \(M\) 纵坐标为 \(2\)，求 \(M\) 与焦点的距离；
\item 若 \(t=-1\)，\(P(1,1)\)，\(Q(1,-1)\)，证明：\(y_Ay_B\) 为常数；
\item 是否存在 \(t\)，使得 \(y_Ay_B=1\) 且 \(y_Py_Q\) 为常数？若存在，求出 \(t\) 的所有可能值；若不存在，请说明理由.
\end{enumerate}
\end{problem}

\begin{solution}
（1） 抛物线 \(y^2=x\) 可写为 \(y^2=4px\)，故 \(p=\dfrac14\)，焦点为 \((\dfrac14,0)\). 当 \(y_0=2\) 时 \(x_0=4\)，由抛物线定义，点到焦点距离为到准线 \(x=-\dfrac14\) 的距离，即 \(4+\dfrac14=\dfrac{17}{4}\).

（2） 设过 \(M\) 与 \(P,Q\) 的直线方程分别写成参数形式，代入 \(x=-1\) 求得 \(A,B\) 的纵坐标. 利用 \(P,Q\) 关于 \(x\) 轴对称，可化简得 \(y_Ay_B\) 与 \(M\) 无关，为常数.

（3） 设直线与抛物线交点参数为 \(u,v\). 用抛物线参数方程 \(x=y^2\) 表示交点和直线与 \(x=t\) 的交点，代入 \(y_Ay_B=1\) 并要求 \(y_Py_Q\) 为常数，解得存在且 \(t=-1\).
\end{solution}

\begin{problem}
已知非空集合 \(A\subseteq\mathbb R\)，函数 \(y=f(x)\) 的定义域为 \(D\). 若对任意 \(t\in A\) 且 \(x\in D\)，不等式 \(f(x)\ge f(x+t)\) 恒成立，则称函数 \(f(x)\) 具有 \(A\) 性质.
\begin{enumerate}
\item 当 \(A=\{-1\}\)，判断 \(f(x)=-x\)，\(g(x)=2^x\) 是否具有 \(A\) 性质；
\item 当 \(A=(0,1)\)，\(f(x)=x+\dfrac1x\)，\(x\in[a,+\infty)\)，若 \(f(x)\) 具有 \(A\) 性质，求 \(a\) 的取值范围；
\item 当 \(A=\{-2,m\}\)，\(m\in\mathbb Z\)，若 \(D\) 为整数集且具有 \(A\) 性质的函数均为常值函数，求所有符合条件的 \(m\) 的值.
\end{enumerate}
\end{problem}

\begin{solution}
（1） 对 \(A=\{-1\}\)，需比较 \(f(x)\) 与 \(f(x-1)\). \(f(x)=-x\) 时，\(-x\ge-(x-1)\) 不成立，故不具有 \(A\) 性质；\(g(x)=2^x\) 时，\(2^x\ge2^{x-1}\) 恒成立，故具有 \(A\) 性质.

（2） 要求对任意 \(t\in(0,1)\)，有 \(f(x)\ge f(x+t)\). 这等价于 \(f\) 在 \([a,+\infty)\) 上单调不增. 因 \(f'(x)=1-\dfrac1{x^2}\)，结合定义域 \(x>0\) 或 \(x<0\) 的限制，可得 \(a\) 的取值范围为 \((-\infty,-1]\).

（3） 在整数集上，具有 \(A\) 性质意味着沿步长 \(-2\) 与 \(m\) 的方向函数值形成单调链. 要所有这样的函数均为常值函数，需由步长 \(-2\) 与 \(m\) 生成整个整数格且能形成双向可达，故 \(\gcd(2,|m|)=1\). 因此 \(m\) 为奇数.
\end{solution}
